\section{Limitations}
\label{sec:limitations}

Chiyoda is currently a stylized research simulator, not an operational
evacuation planner. The station layouts are simplified grids unless imported
from richer geometry. Hazard physics are intentionally lightweight and should
not be read as validated CBRN dispersion models. Physiological impairment uses
piecewise exposure curves that are useful for comparative simulation but not
medical prediction.

The behavioral model is also a hypothesis generator rather than a calibrated
description of a specific crowd. Agents use simplified BDI-style intentions,
local gossip, panic/anxiety states, and familiarity cohorts. These mechanisms
are valuable because they make assumptions inspectable, but the paper must not
claim real-world predictive validity without calibration against evacuation
drills, incident records, controlled VR experiments, or pedestrian trajectory
data.

Finally, entropy is not equivalent to safety. The point of the paper is
precisely that uncertainty reduction can be harmful when it produces synchronized
movement. The proposed information-safety metrics are therefore comparative
diagnostics, not universal welfare functions.

A natural extension is adaptive message generation, including learned or
language-model-mediated guidance. That extension should not be evaluated by
message fluency alone. Under this paper's framing, any generated message is a
safety-control action: it must be replayable, validated against the simulated
state, prevented from inventing unavailable exits or stale hazard claims, and
measured by downstream exposure, queueing, exit balance, and evacuation
outcomes. Without those controls, richer language could make a policy more
persuasive while also amplifying harmful convergence.

\subsection{LLM Extension Protocol}
\label{subsec:llm-extension-protocol}

The repository now includes an optional LLM-guidance policy scaffold, but it is
not part of the completed empirical claims in this paper. Its purpose is to
support a controlled follow-on question: whether generated evacuation messages
add safety value beyond deterministic static, global, entropy-targeted, and
bottleneck-aware baselines. The protocol deliberately treats an LLM as a
bounded message proposer. The simulator exposes a structured state summary,
requires a structured message containing recommended exits, avoided exits,
hazard references, confidence, and optional abstention, and validates the
proposal before it can affect agent beliefs.

This design keeps the extension compatible with the core contribution. A live
API call is not allowed to be the experimental evidence by itself. Paper runs
must use cached or deterministic replay, record the prompt-relevant state and
model metadata, and report rejected messages separately from accepted
interventions. The correct comparison is therefore not whether LLM text sounds
more natural, but whether a validated generated-message policy improves ISE or
reduces HCI under the same information budget and target-selection constraints
as deterministic baselines.

\subsection{Medium LLM Extension Result}
\label{subsec:llm-medium-result}

A first medium-scale LLM extension study applies this protocol to 80 runs:
eight policy variants across ten random seeds. The study includes the
deterministic baselines, a deterministic template generator, three live OpenAI
prompt/validator ablations, and a replay-only verification variant. The live
LLM policies make only eight intervention attempts per run, compared with
32--72 interventions for the deterministic communication policies.

\begin{table}[t]
\centering
\small
\resizebox{\columnwidth}{!}{%
\begin{tabular}{lrrrr}
\toprule
Policy & Evacuated & ISE & HCI & Recipients \\
\midrule
Static beacon & 11.0 & 0.0145 & 6.81 & 503.7 \\
Entropy targeted & 9.2 & 0.0079 & 8.10 & 1064.3 \\
Bottleneck avoidance & 9.4 & 0.0124 & 7.55 & 585.7 \\
Template safety & 9.0 & 0.0155 & 7.86 & 446.6 \\
OpenAI state-only & 9.0 & 0.0414 & 8.39 & 132.3 \\
OpenAI safety-strict & 9.9 & 0.0491 & 7.84 & 133.8 \\
OpenAI entropy-lenient & 10.0 & 0.0429 & 8.76 & 138.8 \\
Replay safety-strict & 9.9 & 0.0491 & 7.84 & 133.8 \\
\bottomrule
\end{tabular}
}
\caption{Medium LLM extension summary over ten seeds. LLM guidance improves
information-safety efficiency mainly by operating under a much smaller
intervention budget, but it does not yet reduce harmful convergence relative
to the strongest conservative baseline.}
\label{tab:llm-medium-extension}
\end{table}

The result is useful but deliberately bounded. The live OpenAI variants
achieve higher information-safety efficiency than the deterministic policies
under this intervention budget, and replay exactly matches the safety-strict
live variant. However, the harmful-convergence index remains high. The
extension therefore supports a cautious paper claim: validated generated
messages can be made replayable and efficient, but richer adaptive language
does not by itself solve crowd synchronization risk. The next experiment
should test whether LLM guidance can reduce HCI under controlled target
selection, hazard severity, and population familiarity regimes.
